\chapter{Conclusiones y trabajo futuro}
\label{cap:conclusiones}

% =====================================================================================
% ⚠️ EXIGENCIA EXPLICITA DE LA NORMA (Instrucciones seccion 2.7):
%    «cada objetivo del trabajo se enlazara con una conclusion». Por eso §9.1 recorre los
%    cinco especificos UNO A UNO y cierra con el general. NO fundirlos en un parrafo.
%
% ⚠️ EXTENSION FIJADA POR EL AUTOR (sep-2026): DOS paginas. Esta AJUSTADO al limite: si se
%    añade una linea, §9.4 se cae a una tercera pagina. Recortar antes de añadir.
%
% ⚠️ NINGUNA CIFRA NUEVA. Todas las de este capitulo estan medidas y redactadas en el
%    capitulo~\ref{cap:resultados}. Si aparece un numero que no este alli, esta inventado.
% =====================================================================================

\section{Cumplimiento de los objetivos}
\label{sec:cumplimiento}

Los cinco objetivos específicos del apartado~\ref{sec:objetivos_especificos} se cierran así:

\begin{enumerate}
    \item \textbf{Comprensión de los fundamentos de la computación cuántica.} Cumplido. El
          Capítulo~\ref{cap:fundamentos} fija lo imprescindible: qué es un qubit, qué mide un
          observable y por qué un circuito es un grafo dirigido.

    \item \textbf{Construcción del conjunto de datos.} Cumplido. Veinte mil circuitos de cinco
          a quince qubits, etiquetados con su degradación exacta sobre la calibración real de
          un procesador de IBM durante cuarenta y dos días, y particionados en tres ejes de
          generalización separables (Capítulo~\ref{cap:planteamiento}).

    \item \textbf{Planteamiento y análisis de la variable objetivo.} Cumplido, y con el
          hallazgo que reorientó el trabajo: la degradación con signo \textbf{no es predecible
          antes de ejecutar}. La reformulación sobre el
          \textbf{factor de supervivencia} sí es aprendible (apartado~\ref{sec:reformulacion}).

    \item \textbf{Realización de la comparativa.} Cumplido, con un matiz sobre el `modelo
          óptimo' que se pedía conseguir: \textbf{no existe uno estrictamente mejor}. Random
          Forest gana en la métrica que decide y es el único que nunca empeora el observable al
          corregirlo, pero las redes de grafos lo superan por márgenes grandes en dos
          situaciones (Capítulo~\ref{cap:resultados}).

    \item \textbf{Cuantificación del aporte de la estructura.} Cumplido, y la respuesta es en
          buena parte negativa. El grafo solo gana donde hay que salirse de lo visto: la media
          de $X$ al extrapolar a catorce y quince qubits (apartado~\ref{sec:res_tamano}) y algunos poco circuitos en observables concretos(apartado~\ref{sec:res_tipo}). En el régimen conocido,
          treinta y cuatro columnas agregadas bastan.
\end{enumerate}

El \textbf{objetivo general} (apartado~\ref{sec:objetivo_general}) se cumple. La
\textbf{hipótesis nula} se refuta en su segunda mitad (los cinco modelos baten con holgura la
línea base del factor constante) y solo en parte en la primera: hay diferencias medibles entre
modelos, pero ninguno domina los tres ejes a la vez.

\section{Aportaciones}
\label{sec:aportaciones}

\begin{enumerate}
    \item \textbf{Un conjunto de datos que no existía.} Su valor no está en el tamaño, sino en
          los tres ejes, que permiten medir por separado tres formas distintas de fallar fuera
          de distribución.

    \item \textbf{La demostración medida de que el objetivo natural del problema no es
          predecible antes de ejecutar}, con su explicación y su reformulación.

    \item \textbf{Una comparativa que separa dos preguntas que suelen confundirse}: cuánto
          acierta un modelo sobre el factor, y cuánto mejora el observable al aplicarlo. La
          segunda no se sigue de la primera, y aquí se mide en lugar de suponerse
          (apartado~\ref{sec:res_mitigacion}).
\end{enumerate}

\section{Limitaciones}
\label{sec:limitaciones_finales}
% NO minimizarlas: declararlas hace creible el resto.

\begin{itemize}
    \item El modelo de ruido \textbf{no incluye crosstalk}: las conclusiones son válidas dentro
          de ese modelo.
    \item No hay validación en \textbf{hardware real}: el plan abierto de IBM da $10$ minutos
          de QPU al mes \cite{ibm2026productos}.
    \item El \textit{ground truth} exige simulación exacta, lo que \textbf{acota el tamaño} de
          los circuitos a quince qubits.
    \item El transpilador de Qiskit \textbf{no es determinista}, así que la reproducibilidad
          del conjunto es \textbf{del plan de generación}.
    \item El ruido se modela como factor multiplicativo puro, sin término afín
          (Anexo~\ref{anexo:factor_negativo}).
\end{itemize}

\section{Trabajo futuro}
\label{sec:futuro}
% Fuente: TFM-Quantum/IDEAS_FUTURAS.md

\begin{itemize}
    \item \textbf{Validación en hardware real}: la única limitación que no se resuelve con más
          cómputo clásico.
    \item \textbf{Incorporar el crosstalk} al modelo de ruido y volver a medir los tres ejes.
    \item \textbf{Modelar el término afín}, para los circuitos en los que el factor cambia de signo.
    \item \textbf{Ampliar la batería de observables}, hoy limitada a cinco.
\end{itemize}
